\chapter{Condições: Tomando Decisões}

\begin{objetivos}
\begin{itemize}
\item Entender o que é uma condição;
\item Usar \textbf{se} e \textbf{se--senão} para o programa decidir;
\item Combinar várias decisões, uma dentro da outra.
\end{itemize}
\end{objetivos}

\section{Pergunta com resposta sim ou não}

Pense nesta frase: \emph{``Se estiver chovendo, leve o guarda-chuva.''} Você só leva o guarda-chuva \textbf{se} a condição for verdadeira. O computador também decide assim!

\begin{vocabulario}[Vocabulário: condição]
\textbf{Condição} é uma pergunta que só tem duas respostas: \textbf{verdadeiro} (sim) ou \textbf{falso} (não). Exemplo: ``o gatinho está tocando na borda?''
\end{vocabulario}

\begin{pensaassim}[Pense assim: o semáforo]
Você chega na rua e olha o semáforo. \emph{Se} estiver verde, atravessa; \emph{senão}, espera. Você \textbf{faz uma pergunta} (``está verde?''), recebe \textbf{sim ou não} e \textbf{escolhe o que fazer}. É exatamente isso que um bloco de condição faz no Scratch!
\end{pensaassim}

\section{O bloco se}

O bloco \bloco[scControl]{se \booleano{...} então} executa os blocos de \textbf{dentro} dele apenas quando a condição é verdadeira.

\begin{codigo}[Se tocar na borda, vira]
\begin{pilha}
\bk[scEvents]{quando \bandeira\ for clicado}
\sbc[scControl]{sempre}
\bk[scMotion]{mova \entrada{5} passos}
\sbc[scControl]{se \booleano{tocando em \menu{borda}?} então}
\bk[scMotion]{gire \entrada{180} graus}
\sbcend
\sbcend
\end{pilha}
\end{codigo}

\section{O bloco se--senão}

Às vezes queremos \textbf{duas saídas}: uma se for verdadeiro e outra se for falso. É o \bloco[scControl]{se \booleano{...} então \ldots senão}.

\begin{pensaassim}[Pense assim: dia de chuva ou de sol]
\emph{Se} estiver chovendo, você leva o guarda-chuva; \emph{senão}, leva os óculos de sol. Sempre acontece \textbf{uma das duas coisas}, nunca as duas ao mesmo tempo e nunca nenhuma.
\end{pensaassim}

Já o bloco \bloco[scControl]{se} sozinho é como ``se estiver frio, vista um casaco'': se não estiver frio, você simplesmente não faz nada.

\begin{codigo}[Se a tecla espaço estiver apertada\ldots]
\begin{pilha}
\bk[scEvents]{quando \bandeira\ for clicado}
\sbc[scControl]{sempre}
\sbc[scControl]{se \booleano{tecla \menu{espaço} pressionada?} então}
\bk[scLooks]{diga \entrada{Apertado!}}
\sbcelse{senão}
\bk[scLooks]{diga \entrada{Solto!}}
\sbcend
\sbcend
\end{pilha}
\end{codigo}

\begin{missao}[Missão 11: Gato obediente]
Faça o gatinho andar \textbf{só enquanto} a seta para a direita estiver apertada. Quando soltar, ele fica parado e diz ``Parei''.
\end{missao}

\begin{codigo}[Andar só quando mandar]
\begin{pilha}
\bk[scEvents]{quando \bandeira\ for clicado}
\sbc[scControl]{sempre}
\sbc[scControl]{se \booleano{tecla \menu{seta direita} pressionada?} então}
\bk[scMotion]{mude x em \entrada{5}}
\bk[scLooks]{diga \entrada{} }
\sbcelse{senão}
\bk[scLooks]{diga \entrada{Parei}}
\sbcend
\sbcend
\end{pilha}
\end{codigo}

\begin{dica}
Quando o programa tem um \bloco[scControl]{sempre}, ele pergunta a mesma coisa várias vezes por segundo. É assim que o computador ``percebe'' mudanças enquanto o jogo roda! É como um \textbf{goleiro}: ele fica de olho na bola o tempo todo, e \emph{se} ela vem, ele pula.
\end{dica}

\section{Decisões dentro de decisões}

Um \bloco[scControl]{se} pode ficar dentro de outro. A primeira condição filtra, a segunda decide.

\begin{pensaassim}[Pense assim: a porta e a chave]
\emph{Se} você chegou em casa, \emph{então} \emph{se} tiver a chave, abre a porta. Só faz sentido perguntar pela chave depois de ter chegado em casa. Um \bloco[scControl]{se} dentro do outro é uma pergunta que só é feita quando a primeira resposta foi ``sim''.
\end{pensaassim}

\begin{codigo}[Dois ses encaixados]
\begin{pilha}
\bk[scEvents]{quando \bandeira\ for clicado}
\sbc[scControl]{sempre}
\sbc[scControl]{se \booleano{tocando em \menu{borda}?} então}
\sbc[scControl]{se \booleano{tecla \menu{espaço} pressionada?} então}
\bk[scLooks]{diga \entrada{Pulei a parede!}}
\sbcend
\sbcend
\sbcend
\end{pilha}
\end{codigo}

\begin{cuidado}
O \bloco[scControl]{se} de dentro só é verificado quando o de fora for verdadeiro. Se você inverter a ordem neste exemplo, muda qual pergunta vem primeiro. A fala ainda exige as duas condições.
\end{cuidado}

\section{Hora de praticar}

\begin{exercicio}[Guarda-chuva]
Faça o gatinho dizer ``Leve o guarda-chuva!'' se estiver tocando na cor azul (``chuva''), e ``Dia de sol!'' se não estiver.
\end{exercicio}

\begin{solucao}
\begin{pilha}
\bk[scEvents]{quando \bandeira\ for clicado}
\sbc[scControl]{sempre}
\sbc[scControl]{se \booleano{tocando na cor \menu{azul}?} então}
\bk[scLooks]{diga \entrada{Leve o guarda-chuva!}}
\sbcelse{senão}
\bk[scLooks]{diga \entrada{Dia de sol!}}
\sbcend
\sbcend
\end{pilha}
\end{solucao}

\begin{exercicio}[Mouse pressionado ou solto]
Clique na bandeira verde. Enquanto o botão do mouse estiver pressionado, o gatinho deve dizer ``Clique!''. Ao soltar, ele deve dizer ``Esperando\ldots''.
\end{exercicio}

\begin{solucao}
\begin{pilha}
\bk[scEvents]{quando \bandeira\ for clicado}
\sbc[scControl]{sempre}
\sbc[scControl]{se \booleano{mouse pressionado?} então}
\bk[scLooks]{diga \entrada{Clique!}}
\sbcelse{senão}
\bk[scLooks]{diga \entrada{Esperando...}}
\sbcend
\sbcend
\end{pilha}
\end{solucao}

\begin{exercicio}[Ping-pong]
Faça o gatinho andar para a direita e voltar para a esquerda ao tocar na borda, sem ficar de cabeça para baixo.
\end{exercicio}

\begin{solucao}
\begin{pilha}
\bk[scEvents]{quando \bandeira\ for clicado}
\bk[scMotion]{defina o estilo de rotação \menu{esquerda-direita}}
\bk[scMotion]{vá para x: \entrada{0} y: \entrada{0}}
\bk[scMotion]{aponte para a direção \entrada{90}}
\sbc[scControl]{sempre}
\bk[scMotion]{mova \entrada{5} passos}
\bk[scMotion]{se tocar na borda, volte}
\sbcend
\end{pilha}
\end{solucao}

\begin{desafio}[Semáforo]
Crie um semáforo: se a tecla \textbf{V} estiver apertada, diga ``Pode ir'' (verde); senão se a tecla \textbf{A} estiver apertada, diga ``Atenção'' (amarelo); senão diga ``Pare'' (vermelho).
\end{desafio}

\begin{solucao}
Um \bloco[scControl]{se} dentro do \bloco[scControl]{senão}:
\begin{pilha}
\bk[scEvents]{quando \bandeira\ for clicado}
\sbc[scControl]{sempre}
\sbc[scControl]{se \booleano{tecla \menu{v} pressionada?} então}
\bk[scLooks]{diga \entrada{Pode ir}}
\sbcelse{senão}
\sbc[scControl]{se \booleano{tecla \menu{a} pressionada?} então}
\bk[scLooks]{diga \entrada{Atenção}}
\sbcelse{senão}
\bk[scLooks]{diga \entrada{Pare}}
\sbcend
\sbcend
\sbcend
\end{pilha}
\end{solucao}
